% ============================================================================
%  Chapitre 5, partie II : Les mesures
% ============================================================================

% ----------------------------------------------------------------------------
\section{Les mesures}
% ----------------------------------------------------------------------------

\subsection{Le gyromètre}

Le gyromètre que l'on va utiliser est le fameux \textbf{MPU6050} de protocol I2C, 
avec un bon timing et une bonne programmation, on peut avoir une fréquence de mesure de $ f=100 $ Hz.

Le gyromètre mesure la vitesse angulaire $\vect{\omega}$ de la boule
dans son propre repère, à la fréquence $f$, donc aux instants
$t_k = k/f$. Une mesure contient un
biais $\vect{b}$, que l'on suppose constant sur la durée d'une
expérience, et un bruit blanc gaussien $\vect{\nu}_k$ :
\[
    \vect{z}_k = \vect{\omega}(t_k) + \vect{b} + \vect{\nu}_k,
    \qquad
    \vect{\nu}_k \sim \mathcal{N}(0, \sigma^2\Id)
\]

\subsection{La verticale vue par la boule}
Les quantités conservées vont intervenir dans l'estimation, et pour $\vect{L} _G \vect{e}_z $, on a besoin de la matrice de changement de repère $\mat{R}$, donc éventuellement un magnétomètre si jamais on a besoin d'orientation de la sphère. Cependant, en supposant que l'on a très peu de bruit sur le gyromètre, une intégration sera suffisante pour réduire l'instalation d'un capteur. En effet, on a 
\[\vect{L}_G \vect{e}_z = \vect{e}_z ^T (\mat{R} \mat{I}_b \mat{R}^T \vect{\Omega}) = \vect{e}_z ^T \mat{R} \mat{I}_b \mat{R}^T \mat{R} \vect{\omega} = (\vect{e}_z ^T\mat{R})(\mat{I}_b \vect{\omega})
\]
Or
\[
    \vect{\gamma} = \mat{R}^T\vect{e}_z
\]
soit
\[
\vect{L}_G \vect{e}_z = \vect{\gamma} ^ T(\mat{I}_b \vect{\omega})
\]
\textbf{Évolution de $\vect{\gamma}$ : } on part de
$\dot{\mat{R}} = \mat{R}\hatx{\vect{\omega}}$. En transposant, et
comme $\hatx{\vect{\omega}}$ est antisymétrique,
\[
    \dot{\mat{R}}^T = \hatx{\vect{\omega}}^T\mat{R}^T
    = -\hatx{\vect{\omega}}\mat{R}^T
\]
Or $\vect{e}_z$ est fixe, donc
\[
    \dot{\vect{\gamma}} = \dot{\mat{R}}^T\vect{e}_z
    = -\hatx{\vect{\omega}}\vect{\gamma}
    = -\vect{\omega}\wedge\vect{\gamma}
\]
soit
\[
    \dot{\vect{\gamma}} = \vect{\gamma}\wedge\vect{\omega}
\]
\textbf{Intégration : } cette équation ne fait intervenir que
$\vect{\omega}$, que mesure le gyromètre. On peut donc obtenir
$\vect{\gamma}$ en l'intégrant à partir d'une valeur initiale
$\vect{\gamma}_0$. Celle-ci est donnée par l'accéléromètre (aussi sur le MPU6050), la boule
étant au repos, il mesure alors $g\,\vect{\gamma}_0$. Entre deux
mesures, on fait la rotation exacte par la vitesse angulaire moyenne
$\bar{\vect{\omega}}_k = \frac12(\vect{z}_k + \vect{z}_{k+1})$ :
\[
    \vect{\gamma}_{k+1}
    = \exp\bigl(-\Delta t\,\hatx{\bar{\vect{\omega}}_k}\bigr)\,\vect{\gamma}_k,
    \qquad \Delta t = 1/f
\]
où l'exponentielle se calcule par la formule de Rodrigues. Une rotation
conserve la norme : on garde $\|\vect{\gamma}_k\| = 1$ sans
renormaliser.

\textbf{Formule de Rodrigues : } pour $\vect{v} \in \R^3$ non nul, on
note $\phi = \|\vect{v}\|$ et $\vect{n} = \vect{v}/\phi$. Comme
$\hatx{\vect{n}}^3 = -\hatx{\vect{n}}$, la série de l'exponentielle se
regroupe en
\[
    \exp\bigl(\hatx{\vect{v}}\bigr)
    = \Id + \sin\phi\,\hatx{\vect{n}} + (1 - \cos\phi)\,\hatx{\vect{n}}^2
\]
c'est la rotation d'angle $\phi$ autour de $\vect{n}$. Ici,
$\vect{v} = -\Delta t\,\bar{\vect{\omega}}_k$ : avec
$\phi_k = \Delta t\,\|\bar{\vect{\omega}}_k\|$ et
$\vect{n}_k = \bar{\vect{\omega}}_k/\|\bar{\vect{\omega}}_k\|$, et
$\hatx{\vect{n}}^2\vect{x} = (\vect{n}\cdot\vect{x})\,\vect{n} - \vect{x}$,
\[
    \vect{\gamma}_{k+1} = \cos\phi_k\,\vect{\gamma}_k
    - \sin\phi_k\,\vect{n}_k\wedge\vect{\gamma}_k
    + (1 - \cos\phi_k)(\vect{n}_k\cdot\vect{\gamma}_k)\,\vect{n}_k
\]

\begin{remarque}[Traitement des données hors ligne]
    $\vect{\gamma}$ est obtenu en intégrant toutes les mesures du
    gyromètre depuis $t_0$, et son erreur grandit avec le temps. On
    traite donc les données hors ligne plutôt qu'en ligne.
\end{remarque}